% Kebei Jiang — Google Novel AI Lead Methodologist
\documentclass[10pt,letterpaper]{article}
\usepackage[margin=0.75in]{geometry}
\usepackage[T1]{fontenc}
\usepackage[utf8]{inputenc}
\usepackage{lmodern}
\usepackage{microtype}
\usepackage{enumitem}
\usepackage[hidelinks]{hyperref}
\usepackage{xcolor}
\usepackage{titlesec}
\usepackage{tabularx}

\usepackage{enumitem}

\pagestyle{empty}
\setlength{\parindent}{0pt}
\setlength{\parskip}{0pt}
\renewcommand{\familydefault}{\sfdefault}

\definecolor{sectiongray}{HTML}{333333}
\titleformat{\section}{\large\bfseries\color{sectiongray}}{}{0pt}{}[\vspace{-0.35em}\rule{\textwidth}{0.5pt}]
\titlespacing*{\section}{0pt}{7pt}{4pt}
\setlist[itemize]{leftmargin=1.25em,itemsep=1.5pt,topsep=1.5pt,parsep=0pt,partopsep=0pt}

\newcommand{\resumeSubheading}[4]{%
\vspace{1pt}
\begin{tabularx}{\textwidth}{@{}Xr@{}}
\textbf{#1} & #2 \\
\textit{#3} & \textit{#4}
\end{tabularx}
\vspace{-3pt}}

\newcommand{\researchHeading}[2]{%
\vspace{5pt}
\begin{tabularx}{\textwidth}{@{}Xr@{}}
\textbf{#1} & \textit{#2}
\end{tabularx}
\vspace{-10pt}}

\begin{document}

\begin{center}
{\LARGE\bfseries Kebei Jiang}\\[3pt]
Bellevue, WA \textbar{} (225) 200-4286 \textbar{}
\href{mailto:kebei.jiang13@gmail.com}{kebei.jiang13@gmail.com} \textbar{}
\href{https://www.linkedin.com/in/kebeijiang}{linkedin.com/in/kebeijiang}
\end{center}

\section{Core Expertise}
\textbf{Evaluation:} LLM-as-a-judge calibration $\cdot$ Agent evaluation $\cdot$ Synthetic benchmarks $\cdot$ Experiment design

\textbf{AI/ML:} LLMs $\cdot$ VLMs $\cdot$ SFT $\cdot$ RLHF/RLVR $\cdot$ RAG

\textbf{Technical:} Python $\cdot$ PyTorch $\cdot$ Braintrust $\cdot$ Hugging Face $\cdot$ Azure $\cdot$ GCP

\section{Experience}

\resumeSubheading{Microsoft}{Jan 2020 -- Sep 2026}{Senior Applied Scientist}{}
\begin{itemize}
\item Developed LLM and VLM solutions spanning agentic AI, RAG, knowledge mining, and model fine-tuning.
\item Led evaluation and experimentation for LLM-based agents and AI systems.
\item Built end-to-end ML and information-processing systems from data preparation through deployment.
\item Led technical development across scientists, PMs, and customer engineers, including technical decisions, implementation reviews, and customer engagements.
\end{itemize}

\resumeSubheading{Viridien}{Jan 2015 -- Jan 2020}{Applied Scientist}{}
\begin{itemize}
\item Applied statistical and machine-learning methods to large-scale seismic data.
\item Developed CNN-based solutions for seismic interpretation.
\end{itemize}

\section{Selected Research \& Evaluation}

\researchHeading{Specification-to-Calibration: Synthetic Supervision for LLM Judges}{2026}
\textit{Submitted to JUDGe @ NeurIPS 2026}
\vspace{10pt}
\begin{itemize}
\item Developed a method for automating LLM judge calibration: compiled natural-language specifications into coverage-controlled positive/negative examples, using targeted behavioral contrasts to revise rubrics while holding the judge model fixed.
\item Improved position-debiased accuracy by 3.58 $\pm$ 0.58 pp on LLMBar; feedback-conditioned and GEPA-based calibration also improved balanced accuracy on BigCodeBench.
\end{itemize}

\researchHeading{Period-Closing Skill --- Stateful, Multi-Turn Agent Evaluation}{2026}
\begin{itemize}
\item Designed 30+ test cases covering diverse failure modes for an LLM agent that automates period-closing workflows involving due dates, calendars, worker assignment, and task dependencies; reviewed cases with the PM to validate representative scenarios.
\item Built realistic multi-turn test environments with pre-existing state, calendar constraints, dependencies, and workers with similar names; fed simulated user responses back into the agent to test multi-turn behavior.
\item Separated execution-trace quality from final task success and added trace/task verification across repeated runs.
\end{itemize}

\researchHeading {Account Reconciliation Agent --- Synthetic Benchmark Methodology} {2026}
\begin{itemize}
    \item Built an evaluation framework for an LLM agent resolving 440 reconciliation exceptions between ledger and subledger using synthetic accounting workflows with ground-truth resolutions and tool-use paths.
    \item Identified an unrealistic uniform distribution of matching cases, validated the expected distribution with the PM, and recalibrated synthetic data generation.
    \item Designed an adaptive preference-monitoring POC using synthetic data with shifting user-preference distributions: trained a tree-based model on a previous window of exception features and user-selected actions, then evaluated it on the current window and triggered retraining when accuracy fell below a PM-approved 90\% threshold.
\end{itemize}

\researchHeading {VLM Fine-Tuning --- Controlled Diagnostic for Shortcut Learning}{2025}
\begin{itemize}
\item Designed a hierarchical taxonomy spanning hundreds of fashion-item classes and aligned it with GPT-4o's understanding of retail products; generated $\sim$1M synthetic training examples and fine-tuned an open-source VLM with LoRA, SFT, and DPO.
\item Designed a diagnostic to test what the model learned: after updating the names of several classes in the taxonomy, the fine-tuned model failed to adopt the changes, exposing taxonomy memorization.
\item Achieved $\sim$90\% accuracy against GPT-4o ground truths; used the diagnostic to motivate separating domain adaptation from taxonomy alignment through distinct training strategies.
\end{itemize}

\researchHeading{Multimodal RAG --- Financial Insights Copilot}{2024}
\begin{itemize}
\item Built an end-to-end RAG system for content-rich financial PDFs, spanning document parsing, multimodal indexing, retrieval, query decomposition, and knowledge-graph exploration.
\item Identified unnecessary intermediate image storage as a document-processing bottleneck and redesigned the indexing flow around image bounding boxes; reduced parsing time by $\sim$80\%.
\item Systematically evaluated multiple BM25 configurations and reranking models and improved offline retrieval by $\sim$20\%; the RAG work was released as part of Microsoft Financial Insights Copilot.
\end{itemize}

\section{Education}
\resumeSubheading{Louisiana State University}{2007 -- 2014}{PhD, Physics}{Baton Rouge, LA}

\section{Publications}
\begin{itemize}
\item \textbf{Guang Yang, Kebei Jiang, David P. Lundquist, Prerna Singh.} \textit{Specification-to-Calibration: Synthetic Supervision for LLM Judges.} Submitted to JUDGe @ NeurIPS 2026.
\item \textbf{Moschoula Pternea, Prerna Singh, Abir Chakraborty, Yagna Oruganti, Mirco Milletari, Sayli Bapat, Kebei Jiang.} \textit{The RL/LLM Taxonomy Tree: Reviewing Synergies Between Reinforcement Learning and Large Language Models.} Journal of Artificial Intelligence Research, 2024. DOI: 10.1613/jair.1.15960.
\item \textbf{Yu Zeng, Kebei Jiang, Jie Chen.} \textit{Automatic Seismic Salt Interpretation with Deep Convolutional Neural Networks.} Proceedings of the 2019 3rd International Conference on Information System and Data Mining, pp. 16--20.
\end{itemize}

\section{Patent}
\begin{itemize}
\item \textbf{Yinheng Li, Kebei Jiang.} \textit{Concurrent Labeling of Sequences of Words and Individual Words.} U.S. patent application, MS\# 411939-US-NP, 2022.
\end{itemize}

\end{document}
